% SGA 3，第 IX 篇，中文重构。
% 范围：第 8 节全文、篇末书目及编者注 49--55。
% 权威底本：Polo--Gille Exp9-8nov09.pdf，组件本地第 30 页（部分）至第 32 页，
% 完整阅读本第 676 页（部分）至第 678 页，印刷第 74--76 页。
% Jacob Reinhold 的 CC BY 4.0 Markdown 谱系仅用于措辞比对。
% 文字、公式、编号和注释均以 Polo--Gille PDF 为准。

\setcounter{footnote}{48}
\SGARefTarget{sga3:ix:section:8}
\section*{8. 域上乘法型群的子群、商群与扩张\textsuperscript{*}}
\addcontentsline{toc}{section}{8. 域上乘法型群的子群、商群与扩张}
\begingroup
\renewcommand{\thefootnote}{*}
\begin{NoHyper}
\footnotetext{1969 年 7 月增补。本节为事后补入，不依赖第 \SGARefLink{sga3:ix:section:3}{3} 至 \SGARefLink{sga3:ix:section:7}{7} 节。}
\end{NoHyper}
\endgroup

% 源页：Exp9 组件本地第 30 页；完整阅读本第 676 页；印刷第 74 页。
\medskip
\noindent\SGARefTarget{sga3:ix:scholium:8:0}\textbf{附论 8.0.}
\footnote{编者注：本附论为新增。} ---
\textit{设 \(k\) 是域，\(H\) 是仿射 \(k\)-群概形。以 \(k[H]\) 表示其仿射代数，以
\[
  X(H)=\operatorname{Hom}_{k\text{-}\mathrm{gr}}
    (H,\mathbf G_{m,k})
\]
表示 \(H\) 的特征标群。由特征标的线性无关引理，\(X(H)\) 是 \(k[H]\) 的线性无关子集。因此，\(H\) 可对角化，当且仅当 \(X(H)\) 张成 \(k\)-向量空间 \(k[H]\)。}

\medskip
\noindent\SGARefTarget{sga3:ix:proposition:8:1}\textbf{命题 8.1.} ---
\textit{设 \(G\) 是域 \(k\) 上的可对角化群概形（分别地，乘法型群概形）。\footnote{编者注：增添了原文中隐含的“在域 \(k\) 上”。此外，由于未假设 \(G\) 有限型，将“代数群”换成了“群概形”。}则对 \(G\) 的每个子群概形 \(H\)，概形 \(H\) 和 \(G/H\) 均可对角化（分别地，均为乘法型）。}

定义 \SGARefLink{sga3:ix:definition:1:1}{1.1} 立即将断言归结到不带括号部分的情形。利用 VI\textsubscript{B} 的 \SGARefLink{prop:sga3-VIB-11.13}{11.13} 和 \SGARefLink{sga3:viii:theorem:3:1}{VIII 3.1}，经简单的极限过渡，可归结到 \(G\) 在 \(k\) 上有限型的情形。\footnote{编者注：应详细写出这一“极限过渡”。它使用 \SGARefLink{sga3:ix:scholium:8:0}{8.0} 及等式
\[
  G/H=\varprojlim_i G_i/H_i
\]
（参见 VI\textsubscript{B}，\SGARefLink{thm:sga3-VIB-11.17}{11.17} 的证明）。为说明 \(H=\varprojlim_i(H\cap G_i)\)，是否使用了 \(H\) 在 \(G\) 中闭这一事实？利用此事实，可给出直接证明。}由 VI\textsubscript{B} 的 \SGARefLink{thm:sga3-VIB-11.16}{11.16}，\footnote{编者注：VI\textsubscript{B} 的 \SGARefLink{thm:sga3-VIB-11.16}{11.16} 无疑应改写成通常的、更方便的形式。特别地，下文一个 \(f_i\) 就足够。}可在 \(G=D_k(M)\) 的仿射环 \(k(M)\) 中找到有限个非零元素
\SGARefTarget{sga3:ix:eq:8:1:1}
\begin{equation}
  \tag{8.1.1}
  f_i=\sum_m a_{im}m,\qquad a_{im}\in k,
\end{equation}
使得 \(H\) 在任意 \(k\)-代数 \(k'\) 中取值的点，恰为满足
\SGARefTarget{sga3:ix:eq:8:1:2}
\begin{equation}
  \tag{8.1.2}
  \tau_g f_i=\lambda_i(g)f_i,
  \qquad \lambda_i(g)\in k',
\end{equation}
的点 \(g\in G(k')\)，其中 \(\tau_g\) 表示按 \(g\) 的平移。然而
\SGARefTarget{sga3:ix:eq:8:1:3}
\begin{equation}
  \tag{8.1.3}
  \tau_g f_i=\sum_m a_{im}\chi_m(g)m,
\end{equation}
其中 \(\chi_m:G\to\mathbf G_{m,k}\) 是对应于 \(m\) 的特征标。因此 \SGARefLink{sga3:ix:eq:8:1:2}{(8.1.2)} 等价于
\SGARefTarget{sga3:ix:eq:8:1:4}
\begin{equation}
  \tag{8.1.4}
  \chi_m(g)=\chi_{m'}(g)
  \qquad\text{若 }m,m'\in Z_i,
\end{equation}
其中 \(Z_i\) 表示使 \(a_{im}\neq0\) 的 \(m\in M\) 所成的集合。\SGARefLink{sga3:ix:eq:8:1:4}{同一关系} 可写成
\[
  \chi_{m'-m}(g)=1
  \qquad\text{若 }m,m'\in Z_i.
\]

% 源页：Exp9 组件本地第 31 页；完整阅读本第 677 页；印刷第 75 页。
设 \(N\) 为 \(M\) 中由所有元素 \(m'-m\) 生成的子群，其中 \(i\) 变动且 \(m,m'\in Z_i\)。由 \(D_k(M)\simeq G\) 的定义，\(H\) 可等同于 \(D_k(M/N)\)。故 \SGARefLink{sga3:viii:theorem:3:1}{VIII 3.1} 将 \(G/H\) 等同于 \(D_k(N)\)。
\hfill\(\square\)

\medskip
\noindent\SGARefTarget{sga3:ix:proposition:8:2}\textbf{命题 8.2.}
\footnote{编者注：陈述及证明均已扩充。} ---
\textit{设 \(k\) 是域；\(H\) 和 \(K\) 是有限型乘法型 \(k\)-群概形；\(G\) 是嵌入正合列
\SGARefTarget{sga3:ix:eq:8:2:0}
\[
  1\longrightarrow H\longrightarrow G\longrightarrow K\longrightarrow 1
\]
中的 \(k\)-群概形（这蕴涵 \(G\) 在 \(k\) 上有限型）。}

\SGARefTarget{sga3:ix:proposition:8:2:item:i}\textit{(i) 若 \(G\) 交换或 \(K\) 连通，则 \(G\) 是乘法型。}

\SGARefTarget{sga3:ix:proposition:8:2:item:ii}\textit{(ii) 若 \(K\) 和 \(H\) 可对角化，且 \(K\) 是环面，则 \(G\) 可对角化。}

关于 \SGARefLink{sga3:ix:proposition:8:2:item:i}{(i)} 的证明，请读者参阅 \SGARefLink{sga3:xvii:proposition:7:1:1}{XVII 7.1.1}，\SGARefLink{sga3:ix:proposition:8:2:item:i}{8.2(i)} 是其特例。域的情形在 \SGARefLink{sga3:xvii:proposition:7:1:1}{XVII 7.1.1} 证明的 \SGARefLink{sga3:xvii:proposition:7:1:1:item:i}{(i)} 部分中处理，不使用后面各篇报告的结果。

对于 \SGARefLink{sga3:ix:proposition:8:2:item:ii}{(ii)}，现假设 \(K\) 和 \(H\) 可对角化：
\[
  K\simeq D_k(M)
  \qquad\text{且}\qquad
  H\simeq D_k(N),
\]
并假设 \(G\) 为乘法型；当 \(K\) 是环面时，这由 \SGARefLink{sga3:ix:proposition:8:2:item:i}{(i)} 得到。由 \SGARefLink{sga3:x:proposition:1:4}{X 1.4}，\(G\) 在 \(k\) 上等平凡；即存在有限可分扩张 \(k'/k\)，使
\[
  G'=G\times_k k'
\]
可对角化。因此，对某个交换群 \(E\)，有 \(G'=D_{k'}(E)\)，且有正合列
\SGARefTarget{sga3:ix:eq:8:2:1}
\begin{equation}
  \tag{1}
  0\longrightarrow D_{k'}(N)\longrightarrow D_{k'}(E)
    \longrightarrow D_{k'}(M)\longrightarrow 0.
\end{equation}
于是由 \SGARefLink{sga3:viii:theorem:3:1}{VIII 3.1} 和 \SGARefLink{sga3:ix:theorem:3:2}{3.2}，\(M\) 是 \(E\) 的子群，并有正合列
\SGARefTarget{sga3:ix:eq:8:2:2}
\begin{equation}
  \tag{2}
  0\longrightarrow M\longrightarrow E\longrightarrow N\longrightarrow 0.
\end{equation}

给定 \(N\) 被 \(M\) 的扩张 \(E_0\)，考虑可对角化 \(k\)-群 \(G_0=D_k(E_0)\)。扩张
\SGARefTarget{sga3:ix:eq:8:2:3}
\begin{equation}
  \tag{3}
  1\longrightarrow H\longrightarrow G_0\longrightarrow K\longrightarrow 1,
\end{equation}
的自同构所成的 \(k\)-群函子 \(A\)，即
\(\operatorname{Aut}_{k\text{-}\mathrm{gr}}(G_0)\)
中在 \(k\)-预概形 \(T\) 上的点为那些在 \(H_T\) 和 \(K_T\) 上诱导恒等的
\[
  \phi\in\operatorname{Aut}_{T\text{-}\mathrm{gr}}(G_T)
\]
的群子函子，可等同于 \(\operatorname{Hom}_{k\text{-}\mathrm{gr}}(K,H)\)。由 \SGARefLink{sga3:viii:corollary:1:5}{VIII 1.5}，这是取值为
\[
  L=\operatorname{Hom}_{\mathrm{gr}}(N,M).
\]
的常值 \(k\)-群。

由此得到，\(K\) 被 \(H\) 的扩张 \(G\) 中，那些在 \(k\) 的可分闭包 \(k_s\) 上同构于 \SGARefLink{sga3:ix:eq:8:2:3}{(3)} 的扩张，与平展拓扑下常值群 \(L_k\) 的 \(k\)-挠子以相同方式分类。置 \(\Gamma=\operatorname{Gal}(k_s/k)\)。由于 \(L\) 是平凡 \(\Gamma\)-模，这些挠子由
\SGARefTarget{sga3:ix:eq:8:2:4}
\[
  H^1_{\mathrm{\acute et}}(k,L_k)
  =H^1(\Gamma,L)
  =\operatorname{Hom}_{\mathrm{gr.\,top.}}(\Gamma,L).
\]
分类。

% 源页：Exp9 组件本地第 32 页；完整阅读本第 678 页；印刷第 76 页。
若 \(K\) 还是环面，则 \(M\) 无挠，从而 \(L\) 也无挠。因此
\[
  \operatorname{Hom}_{\mathrm{gr.\,top.}}(\Gamma,L)=0,
\]
而 \(K\) 被 \(H\) 的每个扩张本来就可对角化。

\medskip
\noindent\SGARefTarget{sga3:ix:remark:8:3}\textbf{注记 8.3.} ---
\textit{如 \SGARefLink{sga3:ix:proposition:8:2}{8.2} 的证明中已指出，\SGARefLink{sga3:ix:proposition:8:2:item:i}{第一个断言} 已在 \SGARefLink{sga3:xvii:proposition:7:1:1}{XVII 7.1.1} 中对任意基底陈述并证明。利用下文 \SGARefLink{sga3:x:lemma:5-13}{X 5.13}，\SGARefLink{sga3:ix:proposition:8:2:item:ii}{第二个断言} 可用本质上相同的证明推广到整正规基概形——或更一般地，几何单分支基概形——的情形。}

\textit{最后，\SGARefLink{sga3:ix:proposition:8:1}{8.1} 也可视为明显更不平凡的结果 \SGARefLink{sga3:ix:theorem:6:8}{6.8} 的推论，后者也对任意基概形有效。}\footnote{编者注：不过，\SGARefLink{sga3:ix:theorem:6:8}{6.8} 的证明使用了 \SGARefLink{sga3:ix:proposition:8:1}{8.1}。}

\SGARefTarget{sga3:ix:bibliography}
\section*{参考文献}
\addcontentsline{toc}{section}{参考文献}
\footnote{编者注：本篇引用的补充参考文献。}

\noindent \SGARefTarget{sga3:ix:bib:bens}[BEns] N. Bourbaki，\textit{集合论〔Théorie des ensembles〕}，第 I--IV 章，Hermann，1970。
